\section{fastText：用轻量分类器做过滤}

\subsection{bag-of-words 的参数问题}

如果把文本直接做 bag-of-words 分类，参数量会非常大。假设词表大小为 $V$，类别数为 $K$，那么一个朴素线性分类器的权重矩阵就可能接近 $V \times K$。对大词表来说，这很快就变成稀疏且昂贵的系统。

\begin{figure}[H]
\centering
\begin{tikzpicture}[every node/.style={font=\small, align=center}]
\node[draw, rounded corners, fill=red!10, minimum width=3cm, minimum height=1cm] (bo) at (0,0) {bag of words\\$V \times K$};
\node[draw, rounded corners, fill=yellow!14, minimum width=3.2cm, minimum height=1cm] (ft) at (5,0) {fastText\\embedding avg + linear head};
\node[draw, rounded corners, fill=green!14, minimum width=3.2cm, minimum height=1cm] (comp) at (10,0) {parameter reduction\\$H(V+K)$};
\draw[-{Latex[length=3mm]}, thick] (bo) -- (ft);
\draw[-{Latex[length=3mm]}, thick] (ft) -- (comp);
\end{tikzpicture}
\caption{fastText 的核心：把高维稀疏问题压缩成低维表示，再用线性头做分类}
\end{figure}

\begin{knowledgebox}{fastText 为什么适合过滤}
\begin{itemize}
    \item 推理快，适合扫描海量网页。
    \item 训练快，适合快速迭代阈值。
    \item 线性结构足够稳，适合二分类过滤。
\end{itemize}
\end{knowledgebox}

\subsection{词嵌入平均}

fastText 的关键就是把词表示成低维 embedding，然后把一个句子的 embedding 做平均，最后接线性分类器。表面看起来简单，实则把“词表巨大”这个问题变成了“维度可控”的问题。

\begin{figure}[H]
\centering
\begin{tikzpicture}[node distance=0.9cm, every node/.style={font=\small, align=center}]
\node[draw, rounded corners, fill=blue!10, minimum width=1.8cm, minimum height=0.8cm] (w1) {word 1};
\node[draw, rounded corners, fill=blue!10, right=of w1, minimum width=1.8cm, minimum height=0.8cm] (w2) {word 2};
\node[draw, rounded corners, fill=blue!10, right=of w2, minimum width=1.8cm, minimum height=0.8cm] (w3) {word 3};
\node[draw, rounded corners, fill=yellow!14, below=1.2cm of w2, minimum width=3.6cm, minimum height=0.8cm] (avg) {average embedding};
\node[draw, rounded corners, fill=green!14, below=of avg, minimum width=3.6cm, minimum height=0.8cm] (cls) {linear classifier};
\draw[-{Latex[length=3mm]}, thick] (w1) -- (avg);
\draw[-{Latex[length=3mm]}, thick] (w2) -- (avg);
\draw[-{Latex[length=3mm]}, thick] (w3) -- (avg);
\draw[-{Latex[length=3mm]}, thick] (avg) -- (cls);
\end{tikzpicture}
\caption{fastText 的前向路径：词嵌入平均后直接线性分类，没有复杂非线性}
\end{figure}

\subsection{ngram 和 hashing trick}

为了捕捉局部短语信息，fastText 还会加入 bigram、trigram 等 n-gram 特征。问题是 n-gram 维度会爆炸，所以要用 hashing trick 把无限词组压进有限桶里。

\begin{figure}[H]
\centering
\begin{tikzpicture}[every node/.style={font=\small, align=center}]
\node[draw, rounded corners, fill=blue!8, minimum width=2.4cm, minimum height=0.8cm] (b1) at (0,1.6) {the cat};
\node[draw, rounded corners, fill=blue!8, minimum width=2.4cm, minimum height=0.8cm] (b2) at (0,0.4) {cat in};
\node[draw, rounded corners, fill=blue!8, minimum width=2.4cm, minimum height=0.8cm] (b3) at (0,-0.8) {in the};
\node[draw, rounded corners, fill=blue!8, minimum width=2.4cm, minimum height=0.8cm] (b4) at (0,-2.0) {the hat};
\node[draw, rounded corners, fill=yellow!16, minimum width=2.6cm, minimum height=1.0cm] (bin1) at (5,0.8) {bin 1};
\node[draw, rounded corners, fill=yellow!16, minimum width=2.6cm, minimum height=1.0cm] (bin2) at (5,-0.8) {bin 2};
\node[draw, rounded corners, fill=yellow!16, minimum width=2.6cm, minimum height=1.0cm] (bin3) at (5,-2.4) {bin 3};
\draw[-{Latex[length=3mm]}, thick] (b1) -- (bin2);
\draw[-{Latex[length=3mm]}, thick] (b2) -- (bin1);
\draw[-{Latex[length=3mm]}, thick] (b3) -- (bin2);
\draw[-{Latex[length=3mm]}, thick] (b4) -- (bin3);
\node[right=1.6cm of bin2] {hash collisions are tolerated};
\end{tikzpicture}
\caption{hashing trick 的作用：把不受控的 n-gram 空间压缩到固定桶数里}
\end{figure}

\begin{warningbox}{fastText 的现实限制}
它本质上仍是一个线性模型。它能快地做“像不像目标类”的判断，但不会自动理解篇章结构、逻辑严密性或事实正确性。
\end{warningbox}

\subsection{在过滤里的角色}

如果 KenLM 更像“语言门卫”，fastText 就更像“轻量分类器门卫”。前者用语言模型分数过滤，后者直接预测文档是否属于目标类。比如语言识别、数学文本过滤、毒性识别，都可以用这个思路。

\begin{figure}[H]
\centering
\begin{tikzpicture}[every node/.style={font=\small, align=center}]
\node[draw, rounded corners, fill=green!10, minimum width=3.2cm, minimum height=0.85cm] (yes) at (0,0) {positive / keep};
\node[draw, rounded corners, fill=red!10, minimum width=3.2cm, minimum height=0.85cm] (no) at (5,0) {negative / drop};
\node[draw, rounded corners, fill=yellow!14, minimum width=3.8cm, minimum height=0.85cm] (binary) at (10.4,0) {for filtering, often just $K=2$};
\draw[-{Latex[length=3mm]}, thick] (yes) -- (binary);
\draw[-{Latex[length=3mm]}, thick] (no) -- (binary);
\end{tikzpicture}
\caption{在过滤任务中，fastText 经常退化成二分类器：保留还是丢弃}
\end{figure}

\subsection{训练目标与概率校准}

fastText 通常优化的是交叉熵：
\[
\mathcal{L}=-\sum_{(x,y)} \log p_\theta(y\mid x).
\]
如果类别不平衡，单纯看准确率没有意义。过滤场景更关心的是 precision、recall 和阈值后剩余样本的质量。换句话说，\emph{你不是在做考试排名，而是在决定哪些数据值得保留。}

\begin{table}[H]
\centering
\small
\renewcommand{\arraystretch}{1.15}
\begin{tabular}{p{2.3cm}p{4.2cm}p{4.1cm}p{3.0cm}}
\toprule
\textbf{指标} & \textbf{说明} & \textbf{过滤里的作用} & \textbf{注意点} \\
\midrule
accuracy & 整体预测对不对 & 不够敏感 & 类别不平衡时失真 \\
precision & 保留下来的样本有多干净 & 很重要 & 常和阈值联动 \\
recall & 目标样本保住了多少 & 防止过度过滤 & 不能只追高 precision \\
calibration & 概率是否可信 & 用于阈值选择 & 域外常会漂移 \\
\bottomrule
\end{tabular}
\caption{fastText 过滤任务里真正有用的指标，不是单纯 accuracy}
\end{table}

\subsection{一个过滤器的工作方式}

你可以把 fastText 当成一个“快速排序器”。它不会替你理解文本，而是把一堆样本按一个近似相关的分数排好，然后让你在高分区间做保留决策。

\begin{enumerate}
    \item 先准备少量标注数据，定义什么叫正样本、负样本。
    \item 训练一个轻量线性分类器，让它学会粗略区分两类文本。
    \item 把模型跑在海量网页上，得到每个样本的置信度。
    \item 在验证集上挑阈值，决定保留多少数据、牺牲多少召回。
\end{enumerate}

\begin{importantbox}{为什么这个流程很稳}
它把“建模”拆成了两个层次：一个是轻量预测器，另一个是阈值控制器。预测器解决次序问题，阈值控制器解决预算问题。两者分开以后，调参会比把所有东西揉成一个大模型更清楚。
\end{importantbox}

\subsection{为什么小模型就够}

fastText 在这里不是为了“理解文本”，而是为了在海量样本上快速给出一个稳定的秩序。只要排序足够好，很多过滤决策就已经足够用了。再往上堆复杂模型，收益往往不如增加标注、调整阈值和清理输入稳定。

\begin{table}[H]
\centering
\small
\renewcommand{\arraystretch}{1.15}
\begin{tabular}{p{2.5cm}p{4.2cm}p{4.0cm}p{3.0cm}}
\toprule
\textbf{模型复杂度} & \textbf{好处} & \textbf{坏处} & \textbf{在过滤里的建议} \\
\midrule
线性/fastText & 便宜、快、稳定 & 表达能力有限 & 优先考虑 \\
小型 Transformer & 更强上下文建模 & 推理成本高 & 只在高价值子任务用 \\
大模型打分器 & 语义强、泛化强 & 很贵、很慢 & 通常不划算 \\
\bottomrule
\end{tabular}
\caption{过滤任务中，模型越强不一定越合适}
\end{table}

\begin{knowledgebox}{fastText 常见的工程优势}
\begin{itemize}
    \item 推理快，适合扫描海量网页。
    \item 训练快，适合快速迭代阈值。
    \item 线性结构足够稳，适合二分类过滤。
    \item 可以用少量标注快速建立基线。
\end{itemize}
\end{knowledgebox}

\subsection{误差结构怎么读}

在过滤场景里，fastText 的假阳性和假阴性都有明确代价。假阳性意味着把坏样本留进来了，假阴性意味着把好样本扔掉了。前者污染训练集，后者降低数据效率。实际系统通常会在这两者之间做偏向性的取舍。

\begin{table}[H]
\centering
\small
\renewcommand{\arraystretch}{1.15}
\begin{tabular}{p{2.7cm}p{4.1cm}p{4.1cm}p{2.8cm}}
\toprule
\textbf{错误类型} & \textbf{含义} & \textbf{后果} & \textbf{常见应对} \\
\midrule
false positive & 坏数据被保留 & 污染训练集 & 提高阈值 \\
false negative & 好数据被丢掉 & 数据利用率下降 & 降低阈值 / 二阶段过滤 \\
calibration drift & 概率不再可信 & 阈值失效 & 重新校准 \\
\bottomrule
\end{tabular}
\caption{fastText 过滤里的三类典型错误与应对}
\end{table}

\begin{warningbox}{过滤任务里的“好”不是绝对的}
一个分类器如果在测试集上看起来很准，但把你真正想保留的长尾样本大量过滤掉，那它在数据工程里依然可能是坏的。过滤的目标不是分类正确率，而是训练集收益。
\end{warningbox}

